%%%%%%%%%%  scan idx 193  |  folio 180  %%%%%%%%%%
\item
Posons $A=\Gamma(S,\mathcal O_S)$, et notons $\Mod(A)$ la catégorie des $A$-modules, $\mathrm D(A)$ la catégorie dérivée correspondante. Le foncteur
\[
\Gamma_S:\Coh(S)\longrightarrow\Mod(A)
\]
est exact, pleinement fidèle et a pour image essentielle la sous-catégorie additive strictement pleine de $\Mod(A)$ formée des $A$-modules de présentation finie. Il induit un foncteur pleinement fidèle
\[
\mathrm D^{-}(\Coh(S))\longrightarrow\mathrm D(A),
\]
dont l'image essentielle est la sous-catégorie pleine $\mathrm D^{-}(A)_{\coh}$ formée des complexes pseudo-cohérents (relativement à la sous-catégorie des modules libres de type fini).
\end{enumerate}

\emph{Preuve.} Pour la première assertion de a) on applique (2.1.2 a)) pour $\mathcal C'_S=\Coh(S)$, l'hypothèse de (2.1.2) étant vérifiée d'après le théorème B : avec les notations de (2.1.2 a)), si $E$ est cohérent et $G$ localement libre de type fini, alors $\underline{\operatorname{Hom}}(G,E)$ est cohérent, donc la catégorie $\Loclib(S)$ (et a fortiori $\Lib(S)$) est relevable dans $\Coh(S)$. La deuxième assertion de a) résulte de (1.1 a)) (resp.(1.2 a)) ) appliqué pour $\mathcal C_0=\Lib$ (resp. $\Loclib$)\footnote{Lib et Loclib sont localement stables par noyau d'épimorphisme d'après (I 2.15.1).}, la condition (ii) respée de (1.1) étant vérifiée d'après le théorème A et (2.1.1 b)). L'assertion b) résulte de (1.1 c)) appliqué successivement pour $\mathcal C_0=\Lib$ et $\mathcal C_0=\Loclib$. La première assertion de c) découle de (2.1.2 b)(ii)) appliqué pour $\mathcal C'_S=\Coh(S)$. Pour prouver que le foncteur $\Gamma_S:\mathrm D^{-}(\Coh(S))\to\mathrm D(A)$ est pleinement fidèle, il suffit, compte tenu de b), de prouver que le composé
\[
\mathrm K^{-}(\Lib(S))\longrightarrow\mathrm D^{-}(\Coh(S))\xrightarrow{\Gamma_S}\mathrm D(A)
\]
l'est. Mais celui-ci se factorise en
\[
\mathrm K^{-}(\Lib(S))\xrightarrow{(*)}\mathrm K^{-}(\Lib(A))\xrightarrow{(**)}\mathrm D(A),
\]
